%% ma: L12-B15-0027
%% lop: 12
%% chuong: 5
%% bai: 15
%% dang: 12.15.22
%% loai: TN4
%% muc: NB
%% dap_an: "C"
%% nguon: "Ảnh giáo viên gửi 24/10/2026 (M06): Kỳ thi thử đợt IV, Thi thử TN THPT 2027 (Biên soạn: Aimbet), Phần 1 Câu 5"
%% kien_thuc: "Bài 14 (phương trình mặt phẳng qua một điểm, biết vectơ pháp tuyến), Bài 15 (vectơ chỉ phương từ phương trình chính tắc)"
%% trang_thai: chinh-thuc
\begin{ex}
Trong không gian với hệ trục tọa độ $Oxyz$, cho mặt phẳng $(P)$ đi qua điểm $M(2;-1;3)$ và vuông góc với đường thẳng $d\colon\dfrac{x+1}{3}=\dfrac{y-2}{-2}=\dfrac{z+4}{1}$. Mặt phẳng $(P)$ có phương trình là
\choice{$3x+2y+z-7=0$}{$3x+2y+z+7=0$}{\True $3x-2y+z-11=0$}{$3x-2y+z+11=0$}
\loigiai{$(P)\perp d$ nên $(P)$ nhận $\vec u=(3;-2;1)$ làm vectơ pháp tuyến: $3(x-2)-2(y+1)+(z-3)=0\Leftrightarrow3x-2y+z-11=0$.}
\end{ex}
